% ===== PRÉAMBULE CANONIQUE — à coller VERBATIM en tête de CHAQUE .tex =====
\documentclass[11pt,a4paper]{article}
\usepackage[utf8]{inputenc}\usepackage[T1]{fontenc}\usepackage{lmodern}
\usepackage[french]{babel}
\usepackage{amsmath,amssymb,mathtools}\usepackage{icomma}
\usepackage[version=4]{mhchem}          % \ce{...} : équations & formules
\usepackage{siunitx}\sisetup{locale=FR} % \qty{}{}, \si{}, \ang{}
\usepackage{chemfig}                    % \chemfig{...} : structures développées
\usepackage{xcolor}\usepackage{colortbl}
\usepackage{tikz}\usepackage{pgfplots}\pgfplotsset{compat=1.18}
\usepackage[most]{tcolorbox}\usepackage{enumitem}\usepackage{array,booktabs}
\usepackage[a4paper,margin=2.2cm]{geometry}
\usetikzlibrary{arrows.meta,calc,angles,quotes,positioning,patterns,backgrounds,shapes.geometric}
\definecolor{primary}{RGB}{222,49,99}\definecolor{primarydark}{RGB}{150,22,55}
\definecolor{defcol}{RGB}{0,114,178}\definecolor{methcol}{RGB}{52,92,148}
\definecolor{excol}{RGB}{0,135,90}\definecolor{attncol}{RGB}{200,40,55}
\tcbset{boxbase/.style={enhanced,breakable,boxrule=0.9pt,arc=2.5pt,
  left=3mm,right=3mm,top=2.3mm,bottom=2.3mm,fonttitle=\bfseries,coltitle=white}}
% --- boîtes TUTORIEL (noms EXACTS, ne pas renommer) ---
\newtcolorbox{cle}[1][]{boxbase,colback=defcol!6,colframe=defcol,
  title={\textsc{À retenir}\ifx\relax#1\relax\relax\else\textmd{ : \itshape #1}\fi}}
\newtcolorbox{methode}[1][]{boxbase,colback=methcol!7,colframe=methcol,sharp corners,
  title={\textsc{Méthode}\ifx\relax#1\relax\relax\else\textmd{ --- \itshape #1}\fi}}
\newtcolorbox{exemplebox}[1][]{boxbase,colback=excol!5,colframe=excol,
  title={\textsc{Exemple}\ifx\relax#1\relax\relax\else\textmd{ : \itshape #1}\fi}}
\newtcolorbox{piege}[1][]{boxbase,colback=attncol!6,colframe=attncol,
  title={\textsc{Piège}\ifx\relax#1\relax\relax\else\textmd{ : \itshape #1}\fi}}
\newtcolorbox{remarque}[1][]{boxbase,colback=black!4,colframe=black!45,
  title={\textsc{Remarque}\ifx\relax#1\relax\relax\else\textmd{ : \itshape #1}\fi}}
% --- boîtes EXERCICE / CORRIGÉ (noms EXACTS, auto-comptées) ---
\newtcolorbox[auto counter]{exo}[1][]{boxbase,colback=primary!4,colframe=primary,
  title={\textsc{Exercice}~\thetcbcounter\ifx\relax#1\relax\relax\else\textmd{ --- \itshape #1}\fi}}
\newtcolorbox[auto counter]{corrige}[1][]{boxbase,colback=excol!4,colframe=excol,
  title={\textsc{Corrigé}~\thetcbcounter\ifx\relax#1\relax\relax\else\textmd{ --- \itshape #1}\fi}}
\newcommand{\R}{\mathbb{R}}\newcommand{\N}{\mathbb{N}}\newcommand{\Z}{\mathbb{Z}}\newcommand{\ds}{\displaystyle}
% (un \usepackage{fancyhdr}+en-tête est possible ; il est ignoré côté web)
% =========================================================================
\begin{document}

\begin{corrige}[une pression donnée en pascals]
On divise par les facteurs de conversion.
\begin{enumerate}[label=\alph*)]
  \item $p=\dfrac{202650}{101325}=\qty{2}{atm}$.
  \item $p=\dfrac{202650}{100000}=\qty{2.0265}{bar}$.
  \item $p=2\times760=\qty{1520}{mmHg}$ (puisque $p=\qty{2}{atm}$).
\end{enumerate}
\end{corrige}

\begin{corrige}[préparer $pV=nRT$ avec $R=\num{0,082}$]
On vise le triplet $(\si{atm},\si{\litre},\si{\kelvin})$.
\begin{enumerate}[label=\alph*)]
  \item $p=\dfrac{760}{760}=\qty{1}{atm}$.
  \item $V=\dfrac{500}{1000}=\qty{0.5}{\litre}$.
  \item $T=25+273=\qty{298}{\kelvin}$.
\end{enumerate}
Les trois grandeurs sont maintenant cohérentes avec $R=\num{0,082}$ ; on pourrait calculer
$n=\dfrac{pV}{RT}$ sans faute d'unité.
\end{corrige}

\begin{corrige}[préparer $pV=nRT$ avec $R=\num{8,314}$]
On vise cette fois le triplet $(\si{\pascal},\si{\meter\cubed},\si{\kelvin})$.
\begin{enumerate}[label=\alph*)]
  \item $p=2\times\num{e5}=\qty{2e5}{\pascal}$ (car $\qty{1}{bar}=\qty{e5}{\pascal}$).
  \item $V=3\times\num{e-3}=\qty{3e-3}{\meter\cubed}$.
  \item $T=127+273=\qty{400}{\kelvin}$.
\end{enumerate}
\end{corrige}

\begin{corrige}[compléter la fiche de conversions de la pression standard]
Valeurs de référence pour $\qty{1}{atm}$.
\begin{enumerate}[label=\alph*)]
  \item $\qty{1}{atm}=\qty{101325}{\pascal}$.
  \item $\qty{1}{atm}\approx\qty{101.3}{\kilo\pascal}$.
  \item $\qty{1}{atm}=\qty{1.013}{bar}$.
  \item $\qty{1}{atm}=\qty{760}{mmHg}$.
\end{enumerate}
\end{corrige}

\end{document}
